\section{Matrices de momentos y realizaciones amplituhedrales de rango superior}
\label{sec:higher-rank}
La partición del registro dodecafásico permite prolongar la construcción
de rango uno a puntos explícitos de rango superior. Los mismos trece
nodos ordenados proporcionan simultáneamente menores de Plücker positivos
y una forma de momentos definida positiva. Estas dos positividades se
obtienen de una sola lista de intervalos; la segunda prueba que la
composición conserva el rango requerido.

Sea \(d_j=K_j/Q\), donde \(Q=\sum_{j=1}^{12}K_j=6263\), y definamos
\[
 t_0=0,\qquad t_j=\sum_{r=1}^jd_r.
\]
La positividad de los doce bloques da
\(0=t_0<t_1<\cdots<t_{12}=1\). Para cada \(2\le k\le9\), ponemos
\[
 C^{(k)}_{ai}=t_i^a\quad(0\le a<k),\qquad
 Z^{(k)}_{bi}=t_i^b\quad(0\le b<k+4),\qquad 0\le i\le12.
\]
La convención \(t_i^0=1\) incluye el nodo cero. Así, \(C^{(k)}\)
es una matriz de tamaño \(k\times13\), mientras \(Z^{(k)}\) tiene
tamaño \((k+4)\times13\).

\begin{theorem}[Positividad y rango de la realización de momentos]
La matriz \(C^{(k)}\) representa un punto de
\(\Gr_{>0}(k,13)\), y todos los menores máximos ordenados de
\(Z^{(k)}\) son positivos. La matriz
\[
 Y^{(k)}=C^{(k)}(Z^{(k)})^{\mathsf T},\qquad
 Y^{(k)}_{ab}=\sum_{i=0}^{12}t_i^{a+b},
\]
tiene rango \(k\). En particular, para el dato externo positivo
\(Z^{(k)}\) ya construido a partir del registro,
\[
 [Y^{(k)}]\in\mathcal A_{13,k,4}(Z^{(k)}),\qquad 2\le k\le9,
\]
donde
\[
 \mathcal A_{13,k,4}(Z^{(k)})
 =\{[C(Z^{(k)})^{\mathsf T}]:[C]\in\Gr_{\ge0}(k,13)\}.
\]
\end{theorem}
\begin{proof}
Un menor máximo de cualquiera de las dos matrices iniciales tiene la forma
\[
 \det(t_{i_j}^{a})_{a=0}^{r-1}
 =\prod_{p<q}(t_{i_q}-t_{i_p})>0,
\]
para índices crecientes. La desigualdad prueba las dos afirmaciones sobre
Plücker. El bloque inicial \(k\times k\) de \(Y^{(k)}\) es la matriz
de Hankel \(H^{(k)}\). Para \(x\ne0\),
\[
 x^{\mathsf T}H^{(k)}x
 =\sum_{i=0}^{12}\left(\sum_{a=0}^{k-1}x_at_i^a\right)^2>0.
\]
En efecto, un polinomio no nulo de grado inferior a \(k\le9\) no se
anula en trece puntos distintos. Luego \(H^{(k)}\) es definida positiva,
su determinante es estrictamente positivo y \(Y^{(k)}\) tiene rango
\(k\). La pertenencia amplituhedral es entonces la composición de las
dos matrices positivas anteriores, con el rango ya comprobado.
\end{proof}

En rango dos, la prueba adopta la forma particularmente explícita
\[
 \det H^{(2)}=13\sum_i t_i^2-\left(\sum_i t_i\right)^2
 =\sum_{i<j}(t_j-t_i)^2>0.
\]
La separación de los nodos, producida por la positividad del registro,
queda así convertida en coercividad de una forma cuadrática. A cada
registro fijo corresponde un punto determinado para cada rango. La
cobertura completa, las triangulaciones y los residuos demostrados en la
sección anterior mantienen su dominio de rango uno; el presente teorema
añade la colección explícita de puntos de rangos dos a nueve y sus matrices
de momentos. La distinción permite reutilizar cada demostración con su
alcance exacto.
